\subsection{Ontology Learning}
\label{sec:ontology-learning}
\slidesrc{Ontology learning: 1--25; Ontology development: 74}

Building an ontology by hand (Section~\ref{sec:design-process}) is slow and expensive.
\textbf{Ontology learning} uses NLP and machine learning to do part of the work
automatically. The deck is a tutorial by Ícaro Medeiros (CIn--UFPE, 2008).

\subsubsection{What, why and from where}
\slidesrc{Ontology learning: 4--9}

The field has \textbf{too many names for the same subject}: ontology extraction,
emergence, generation, acquisition, discovery, population, enrichment\ldots\ The common
name, after Philipp Cimiano's book \emph{Ontology Learning and Population from Text}
(2006), is \textbf{ontology learning} (OL).

\begin{nbdef}{Ontology learning}
\textbf{Ontology learning} is the set of methods and techniques for
\begin{itemize}[nosep]
  \item building an ontology from scratch;
  \item enriching or adapting an existing ontology;
\end{itemize}
by extracting concepts and relations (from text or other sources) to form an ontology.
OL is a \textbf{semi-automatic} task of information extraction: a human remains in the
loop to validate and correct.
\end{nbdef}

\paragraph{Why?}
Ontology engineering has three problems (Maedche and Staab, 2001):
\begin{itemize}[nosep]
  \item \textbf{time}: can you develop an ontology fast?
  \item \textbf{difficulty}: is it difficult to build an ontology?
  \item \textbf{confidence}: how do you know that you have got the ontology right?
\end{itemize}
OL can help with all three, especially with the \textbf{knowledge acquisition
bottleneck}: experts are few and busy, and writing their knowledge down by hand is the
slowest part of building any knowledge-based system.

\paragraph{Information sources.}
Relevant text (mainly Web documents); Web document schemata (XML, DTD, RDF); databases
on the Web; dictionaries; semi-structured documents; personal wikis, e-mail and file
folders; existing Web ontologies.

\begin{center}
\begin{tikzpicture}[st/.style={draw=nborange, fill=nborange!10, font=\small,
    rounded corners=2pt, inner sep=4pt}, arr/.style={-{Stealth}, very thick, nbteal}]
  \draw[line width=14pt, black!12] (0,0) circle (2.6);
  \node[font=\scriptsize\itshape] at (0,2.6) {Legacy + Application Data};
  \node[font=\scriptsize\itshape] at (0,-2.6) {Legacy + Application Data};
  \node[draw, circle, fill=nbteal!10, font=\scriptsize, align=center, minimum size=1.7cm]
       at (0,0) {Domain\\ Ontology};
  \draw[arr] (150:1.6) arc (150:30:1.6);
  \draw[arr] (-30:1.6) arc (-30:-150:1.6);
  \node[font=\scriptsize\itshape, text=nbteal] at (0,1.25) {ontology learning};
  \node[font=\scriptsize\itshape, text=nbteal] at (0,-1.25) {ontology learning};
  \node[st] at (150:2.6) {Extract};
  \node[st] at (40:2.6) {Prune};
  \node[st] at (-20:2.6) {Refine};
  \node[st] at (-60:2.6) {Apply};
  \node[st] at (210:2.6) {Import / Reuse};
\end{tikzpicture}
\captionof{figure}{The ontology engineering cycle of Maedche and Staab (2001), redrawn:
import/reuse existing ontologies, extract, prune, refine, apply, and start again.
Ontology learning is not only the task of extraction.}
\end{center}

The cycle shows that OL is \textbf{not only extraction}: one also \emph{imports} and
reuses existing ontologies, \emph{prunes} what is irrelevant for the application,
\emph{refines} the result, and \emph{applies} it, which produces new data for the next
round.

\subsubsection{How to learn ontologies}
\slidesrc{Ontology learning: 10--15}

The main families of techniques: natural language processing; dictionary parsing;
statistical analysis; machine learning; hierarchical concept clustering; Formal Concept
Analysis (lattices).

\paragraph{Why text?}
Text is massively available on the Web; relevant texts contain relevant knowledge about a
domain; and when an ontology is learnt from text, the \emph{linguistic} knowledge (how
concepts are named and talked about) remains associated with the ontology (Sintek et
al., 2004).

\begin{nbintuition}[OL as reverse engineering (Buitelaar et al., 2005)]
People share a model of the world (a ``shared world model'') and use it to
\emph{write} texts: newspapers, manuals, Web pages. Ontology learning goes the opposite
way: from the texts back to the model. It is \textbf{reverse engineering}: like
recovering the design of a program from its compiled code, we recover the conceptual
model from the documents it produced. And like reverse engineering, it is only partial:
much of the shared knowledge is never written down because everybody already knows it.
\end{nbintuition}

\begin{center}
\begin{tikzpicture}[lay/.style={draw, fill=nbblue!6, font=\small, minimum height=0.75cm,
    anchor=east}]
  \foreach \t/\w/\y/\ex in {
      Terms/6.1/0/{disease, illness, hospital},
      (Multilingual) Synonyms/5.8/0.75/{\{disease, illness, Krankheit\}},
      Concepts/5.5/1.5/{DISEASE := $\langle$Int, Ext, Lex$\rangle$},
      Taxonomy/5.2/2.25/{is\_a(DOCTOR, PERSON)},
      Relations/4.9/3.0/{cure(dom: DOCTOR, range: DISEASE)},
      Rules \& Axioms/4.6/3.75/{$\forall x, y\,(\mathit{sufferFrom}(x,y) \to \mathit{ill}(x))$}}
  {
    \node[lay, minimum width=\w cm - 1.6cm] at (12,\y) {\t};
    \node[font=\small\itshape, anchor=east] at ({12-(\w-1.6)-0.15},\y) {\ex};
  }
\end{tikzpicture}
\captionof{figure}{The ontology learning layer cake (Buitelaar et al., 2005), redrawn.
Each layer is a learning task that builds on the layers below it; on the left, an
example of the output of each task.}
\label{fig:layer-cake}
\end{center}

\begin{nbfact}{the layer cake}
Ontology learning from text is organised in six layers, from the bottom:
\begin{enumerate}[nosep]
  \item \textbf{terms}: the relevant words and phrases of the domain (\emph{disease},
        \emph{illness}, \emph{hospital});
  \item \textbf{(multilingual) synonyms}: groups of terms with the same meaning, also
        across languages (\emph{disease}, \emph{illness}, German \emph{Krankheit});
  \item \textbf{concepts}: a concept is defined by its \textbf{int}ension, its
        \textbf{ext}ension and its \textbf{lex}ical realisations (see below);
  \item \textbf{taxonomy}: the is-a hierarchy between concepts (a doctor is a person);
  \item \textbf{relations}: other relations between concepts, with domain and range
        (doctors cure diseases);
  \item \textbf{rules and axioms}: general logical statements (whoever suffers from
        something is ill).
\end{enumerate}
Lower layers are easier to learn automatically and more reliable; upper layers are
harder and need more human validation.
\end{nbfact}

\subsubsection{Term extraction}
\slidesrc{Ontology learning: 16--18}

\textbf{Linguistic methods} reuse the syntactic tools of the course:
\begin{itemize}[nosep]
  \item \textbf{part-of-speech tagging} (Section~\ref{sec:pos-tagging}) to identify the
        syntactic class: a noun suggests a class, a verb a relation;
  \item \textbf{stemming} (Section~\ref{sec:stemming-lemmatization}): \emph{formalize},
        \emph{formalization}, \emph{formalized}, \emph{formalizing} are one term;
  \item \textbf{head-modifier analysis}: in \emph{fast car} the head is \emph{car} and
        \emph{fast} a modifier; in \emph{the hood of the car}, \emph{hood} is a part of
        \emph{car};
  \item \textbf{grammatical function analysis}: ``John played football in the garden''
        $\to$ \emph{play(John, football)}: the subject and the object of a verb become
        the arguments of a relation.
\end{itemize}

\textbf{Statistical methods}:
\begin{itemize}[nosep]
  \item \textbf{term weighting} with TF-IDF (Section~\ref{sec:stop-words-tfidf});
  \item \textbf{co-occurrence analysis} (a common method in text mining);
  \item \textbf{comparison of frequencies between a domain corpus and a general corpus}:
        a word much more frequent in medical texts than in newspapers (\emph{tumour}) is
        likely a medical term; a word equally frequent in both (\emph{table}) is not.
\end{itemize}

\textbf{Hybrid methods}: linguistic rules extract term candidates (e.g.\ the POS pattern
``adjective* noun+'' finds \emph{heart disease}, \emph{chronic kidney disease}), then
statistics filter them (before or after).

\subsubsection{Synonym extraction}
\slidesrc{Ontology learning: 19--20}

Two approaches: \textbf{extending WordNet} (\emph{term classification}: attach a new
term to the right synset of WordNet, Section~\ref{sec:wordnet}) and
\textbf{co-occurrence between terms} (\emph{term clustering}: terms with similar
contexts are grouped together, the distributional hypothesis of
Section~\ref{sec:distributional-hypothesis}).

\subsubsection{Concept extraction}
\slidesrc{Ontology learning: 21--25}

A term may indicate a concept, if we define its:
\begin{itemize}[nosep]
  \item \textbf{intension}: an (in)formal definition of the objects the concept
        describes. E.g.\ ``a disease is an impairment of health or a condition of
        abnormal functioning'';
  \item \textbf{extension}: the set of objects described by the concept. E.g.\ cancer,
        heart disease;
  \item \textbf{lexical realisations}: the term itself and its (multilingual)
        synonyms.
\end{itemize}

\begin{nbintuition}[intension vs.\ extension]
The \emph{intension} is the \emph{recipe} (``what it takes to be a disease''); the
\emph{extension} is the \emph{list} (``all the diseases''). The concept ``even number
between 1 and 9'' has intension ``divisible by 2, between 1 and 9'' and extension
$\{2, 4, 6, 8\}$. In an ontology, the class definition is the intension, and the
instances are the extension.
\end{nbintuition}

\paragraph{Learning the intension.}
An \emph{informal} definition is a shallow definition as used in WordNet (a gloss):
find the appropriate WordNet concept for a term and the appropriate conceptual relations
(Navigli and Velardi, 2004). A \emph{formal} definition gives formal constraints for
class membership, e.g.\ with Formal Concept Analysis.

\paragraph{Learning the extension: ontology population.}
Extracting the instances of a concept from text is \textbf{ontology population}. It is
related to knowledge markup and tag suggestion (semantic metadata), and it uses
\textbf{named entity recognition} (Section~\ref{sec:named-entity-recognition}):
``John is a football player'' $\to$ \emph{John} (a Person) is an instance of
\emph{Football Player}. Instances can be
\begin{itemize}[nosep]
  \item names for objects: persons, organisations, countries, cities;
  \item event instances: a football match, with its teams, players, officials, \ldots
\end{itemize}

\begin{nbwarning}[what the shortened deck leaves out]
The outline of the tutorial continues with taxonomy, relation and axiom learning,
ontology learning from folksonomies (collaborative tags), tools and conclusions. These
parts are not in the shortened deck shown in class; the layer cake of
Figure~\ref{fig:layer-cake} gives what these tasks produce.
\end{nbwarning}

\begin{nbsummary}[ontology learning]
\begin{itemize}[nosep]
  \item OL (a.k.a.\ extraction, generation, acquisition, population, enrichment):
        semi-automatic building or enrichment of ontologies; fights the knowledge
        acquisition bottleneck (time, difficulty, confidence).
  \item Sources: Web text, schemata, databases, dictionaries, semi-structured documents,
        existing ontologies. OE cycle: import/reuse, extract, prune, refine, apply.
  \item Methods: NLP, dictionary parsing, statistics, ML, clustering, Formal Concept
        Analysis. OL from text = reverse engineering of the shared world model.
  \item Layer cake: terms $\to$ synonyms $\to$ concepts $\to$ taxonomy $\to$ relations
        $\to$ rules \& axioms.
  \item Terms: POS, stemming, head-modifier, grammatical functions; TF-IDF,
        co-occurrence, domain vs.\ general corpus; hybrid.
  \item Concept = intension + extension + lexical realisations; extension learning =
        ontology population, with NER.
\end{itemize}
\end{nbsummary}
